\documentclass[10pt]{article}

\usepackage[utf8]{inputenc}
\usepackage[T1]{fontenc}
\usepackage[spanish]{babel}
\usepackage{amsmath,amssymb}
\usepackage[margin=1.45cm]{geometry}
\usepackage{xcolor}
\usepackage{microtype}
\usepackage{enumitem}
\pagestyle{empty}
\setlength{\parindent}{0pt}
\setlength{\parskip}{3pt}
\setlist[enumerate]{leftmargin=*,itemsep=2pt,topsep=2pt}
\newcommand{\Rb}{R_B}

\begin{document}

\begin{center}
{\large\bfseries Derivación de la conclusión principal del modelo}\\[-1pt]
{\small ¿Cuándo el acceso a financiamiento alternativo aumenta la inversión?}
\end{center}

\textbf{1. Problema de la empresa.}
La firma tiene fondos internos \(n\), puede pedir crédito bancario \(b\) hasta
un límite \(\bar b\), y puede obtener financiamiento alternativo \(m\geq 0\).
El banco cobra \(\Rb\) y la fuente alternativa cobra \(\Rb+\tau\), donde
\(\tau\geq 0\) es el costo adicional de acceso. La firma elige
\[
 \max_{k,b,m}\;
 \beta F(k)-\Rb b-(\Rb+\tau)m
 \quad\text{s.a.}\quad
 k=n+b+m,\qquad 0\leq b\leq\bar b,\quad m\geq0,
\]
con \(F'(k)>0\) y \(F''(k)<0\). Como \(\Rb\leq\Rb+\tau\), la firma usa primero
el crédito bancario y solo después la fuente alternativa.

\textbf{2. Capital disponible al agotar el banco.}
Si la firma usa toda su línea bancaria, alcanza
\[
 \bar k\equiv n+\bar b.
\]
Nos concentramos en una firma que todavía querría invertir más si pudiera
financiar la siguiente unidad al costo bancario:
\[
 \beta F'(\bar k)>\Rb.
\]
Por tanto, la pregunta relevante es si el rendimiento de una unidad adicional,
\(\beta F'(\bar k)\), cubre su costo alternativo, \(\Rb+\tau\).

\textbf{3. La comparación marginal genera el umbral.}
Inmediatamente a la derecha de \(\bar k\), la pendiente del beneficio es
\[
 \Pi'_+(\bar k)
 =\beta F'(\bar k)-(\Rb+\tau).
\]
Definimos el costo de acceso máximo que la firma está dispuesta a pagar:
\[
 \boxed{\tau^*\equiv\beta F'(\bar k)-\Rb>0}.
\]
Sustituyendo esta definición,
\[
 \boxed{\Pi'_+(\bar k)=\tau^*-\tau}.
\]
Esta igualdad resume el mecanismo completo:

\begin{enumerate}
 \item Si \(\tau\geq\tau^*\), entonces
 \(\Pi'_+(\bar k)\leq0\). Continuar después del límite no es rentable.
 Como el beneficio es cóncavo, la firma se detiene en
 \[
   \boxed{k^*=\bar k=n+\bar b.}
 \]
 En este régimen, ampliar la línea bancaria sí aumenta la inversión.

 \item Si \(0\leq\tau<\tau^*\), entonces
 \(\Pi'_+(\bar k)>0\). La firma usa financiamiento alternativo e invierte
 más que \(\bar k\), hasta que
 \[
   \boxed{\beta F'(k^*)=\Rb+\tau,\qquad k^*>\bar k.}
 \]
 Aquí una línea bancaria mayor sustituye financiamiento alternativo, mientras
 que una reducción de \(\tau\) eleva la inversión total.
\end{enumerate}

\textbf{4. Magnitud en el caso cuadrático.}
Si \(F(k)=ak-\frac{c}{2}k^2\), con \(c>0\), entonces
\[
 \tau^*=\beta[a-c(n+\bar b)]-\Rb,\qquad
 k^*-\bar k=
 \boxed{\max\left\{0,\frac{\tau^*-\tau}{\beta c}\right\}}.
\]
Así, no basta con que exista financiamiento alternativo: su costo determina
si el límite bancario restringe la inversión y, cuando no la restringe,
cuánto capital adicional instala la firma.

\end{document}
